\hypertarget{python-2.0-to-2.1-comprehensions-nested-scopes-cycle-detecting-gc}{%
\chapter{Python 2.0 to 2.1: Comprehensions, Nested Scopes, \&
Cycle-Detecting
GC}\label{python-2.0-to-2.1-comprehensions-nested-scopes-cycle-detecting-gc}}

\hypertarget{python-2.0-list-comprehensions-scope-leakage}{%
\subsection{3.1 Python 2.0 List Comprehensions \& Scope
Leakage}\label{python-2.0-list-comprehensions-scope-leakage}}

Introduced in Python 2.0, list comprehensions provided a concise syntax
to create lists from sequences. However, their early implementation in
the compiler had a major design flaw: they leaked loop variables into
the surrounding scope.

\hypertarget{inline-loop-compilation}{%
\subsubsection{1. Inline Loop
Compilation}\label{inline-loop-compilation}}

In Python 2.x, a list comprehension did not create a new stack frame or
lexical scope. Instead, it was compiled as an inline loop inside the
host function. Consider the following Python 2.x list comprehension:

\begin{lstlisting}[language=Python]
# Python 2.x Behavior
x = 999
data = [x for x in range(5)]
print(x)  # Outputs 4! The variable x was overwritten.
\end{lstlisting}

\hypertarget{bytecode-analysis}{%
\subsubsection{2. Bytecode Analysis}\label{bytecode-analysis}}

If we trace the bytecode of this list comprehension under Python 2.x, we
see the compiler emits direct modifications to the local frame variables
array:

\begin{lstlisting}
  2           0 BUILD_LIST               0
              3 LOAD_GLOBAL              0 (range)
              6 LOAD_CONST               1 (5)
              9 CALL_FUNCTION            1
             12 GET_ITER
        >>   13 FOR_ITER                12 (to 28)
             16 STORE_FAST               0 (x)       # Direct write to host local variable x!
             19 LOAD_FAST                0 (x)
             22 LIST_APPEND              2
             25 JUMP_ABSOLUTE           13
        >>   28 STORE_FAST               1 (data)
\end{lstlisting}

Because it uses \passthrough{\lstinline!STORE\_FAST 0 (x)!}, the loop
index directly overwrites the host function's local variable
\passthrough{\lstinline!x!}.

\hypertarget{modern-python-3-isolation}{%
\subsubsection{3. Modern Python 3
Isolation}\label{modern-python-3-isolation}}

To resolve this leakage, Python 3.0 redefined list, set, and dictionary
comprehensions to execute inside a compiler-generated nested function
scope. This creates a separate \passthrough{\lstinline!PyFrameObject!}
during evaluation, protecting the host function's namespace from
leakage.

\begin{center}\rule{0.5\linewidth}{0.5pt}\end{center}

\hypertarget{pep-227-nested-scopes-and-legb-name-resolution}{%
\subsection{3.2 PEP 227 Nested Scopes and LEGB Name
Resolution}\label{pep-227-nested-scopes-and-legb-name-resolution}}

Before Python 2.1, CPython used a simple three-tier lookup model:
\textbf{LGB (Local, Global, Built-in)}. * \textbf{The Limitation}: If a
function was nested inside another function, the inner function could
not access local variables defined in the outer function.

\begin{lstlisting}[language=Python]
# Pre-Python 2.1 Behavior
def outer():
    x = 10
    def inner():
        return x  # Raised NameError! x was not in local, global, or builtins.
\end{lstlisting}

To access outer variables, developers had to pass them explicitly as
default arguments (e.g., \passthrough{\lstinline!def inner(x=x):!}).

\hypertarget{the-legb-resolution-engine-pep-227}{%
\subsubsection{1. The LEGB Resolution Engine (PEP
227)}\label{the-legb-resolution-engine-pep-227}}

Introduced as an optional feature in Python 2.1 (enabled via
\passthrough{\lstinline!from \_\_future\_\_ import nested\_scopes!}) and
made standard in Python 2.2, PEP 227 added the \textbf{Enclosing} scope
to name resolution, establishing the modern \textbf{LEGB} lookup
hierarchy:

\begin{lstlisting}
                  [Name Resolution Lookup Path]
                                |
             1. Local Scope (Array offset / LOAD_FAST)
                                |
             2. Enclosing Scopes (Cell reference / LOAD_DEREF)
                                |
             3. Global Scope (Module dict / LOAD_GLOBAL)
                                |
             4. Built-in Scope (Builtins dict / LOAD_GLOBAL)
\end{lstlisting}

\begin{enumerate}
\def\labelenumi{\arabic{enumi}.}
\tightlist
\item
  \textbf{Local (L)}: Fast array offset reads inside the active
  \passthrough{\lstinline!PyFrameObject!}.
\item
  \textbf{Enclosing (E)}: Searches cells
  (\passthrough{\lstinline!PyCellObject!}) containing enclosed variable
  pointers.
\item
  \textbf{Global (G)}: Searches the module globals dictionary
  (\passthrough{\lstinline!\_\_dict\_\_!}).
\item
  \textbf{Built-in (B)}: Searches the standard built-in namespace
  dictionary (\passthrough{\lstinline!\_\_builtins\_\_!}).
\end{enumerate}

\begin{center}\rule{0.5\linewidth}{0.5pt}\end{center}

\hypertarget{closure-mechanics-and-bytecode-compilation}{%
\subsection{3.3 Closure Mechanics and Bytecode
Compilation}\label{closure-mechanics-and-bytecode-compilation}}

When a function accesses variables from an enclosing scope, CPython
compiles the enclosing variables as \textbf{free variables} and stores
them on the heap inside \passthrough{\lstinline!PyCellObject!} structs.

\hypertarget{bytecode-compilation-trace}{%
\subsubsection{1. Bytecode Compilation
Trace}\label{bytecode-compilation-trace}}

Consider the following nested scope implementation:

\begin{lstlisting}[language=Python]
def parent(x):
    def child():
        return x
    return child
\end{lstlisting}

The CPython compiler processes this nested structure using three
specialized bytecode instructions:
\passthrough{\lstinline!LOAD\_CLOSURE!},
\passthrough{\lstinline!MAKE\_FUNCTION!} (with flags), and
\passthrough{\lstinline!LOAD\_DEREF!}.

Let's trace the bytecode for both the \passthrough{\lstinline!parent!}
and \passthrough{\lstinline!child!} functions:

\textbf{Bytecode for \passthrough{\lstinline!parent!}:}

\begin{lstlisting}
  2           0 LOAD_CLOSURE             0 (x)       /* Load cell reference for cell variable x */
              2 BUILD_TUPLE              1           /* Package cell reference inside a tuple */
              4 LOAD_CONST               1 (<code object child>)
              6 LOAD_CONST               2 ('parent.<locals>.child')
              8 MAKE_FUNCTION            8           /* Flag 8 tells the VM to bind closure cells tuple */
             10 STORE_FAST               1 (child)

  4          12 LOAD_FAST                1 (child)
             14 RETURN_VALUE
\end{lstlisting}

\textbf{Bytecode for \passthrough{\lstinline!child!}:}

\begin{lstlisting}
  3           0 LOAD_DEREF               0 (x)       /* Dereference closure cell and load value */
              2 RETURN_VALUE
\end{lstlisting}

\hypertarget{instruction-walkthrough}{%
\subsubsection{2. Instruction
Walkthrough}\label{instruction-walkthrough}}

\begin{enumerate}
\def\labelenumi{\arabic{enumi}.}
\tightlist
\item
  \textbf{\passthrough{\lstinline!LOAD\_CLOSURE!}}: Pushes the cell
  reference representing the cell variable \passthrough{\lstinline!x!}
  onto the stack.
\item
  \textbf{\passthrough{\lstinline!MAKE\_FUNCTION 8!}}: Creates a new
  function object (\passthrough{\lstinline!PyFunctionObject!}). The flag
  \passthrough{\lstinline!8!} (in modern CPython, this uses
  \passthrough{\lstinline!MAKE\_FUNCTION!} with cell tuples on the
  stack) instructs the VM to pop the cells tuple and assign it to the
  function's \passthrough{\lstinline!\_\_closure\_\_!} attribute.
\item
  \textbf{\passthrough{\lstinline!LOAD\_DEREF!}}: Executed inside
  \passthrough{\lstinline!child!}. The VM bypasses local dictionary or
  local array lookups. It accesses the closure tuple at index 0,
  dereferences the \passthrough{\lstinline!PyCellObject!}, and pushes
  its \passthrough{\lstinline!ob\_ref!} payload onto the evaluation
  stack.
\end{enumerate}

\begin{lstlisting}[language=Python]
import inspect

closure_fn = parent(42)
# Inspect closure cells at runtime
print("Closure cells:", closure_fn.__closure__)
print("Cell value:", closure_fn.__closure__[0].cell_contents)
print("Parent cell vars:", parent.__code__.co_cellvars)
print("Child free vars:", closure_fn.__code__.co_freevars)
\end{lstlisting}

\begin{center}\rule{0.5\linewidth}{0.5pt}\end{center}

\hypertarget{python-2.0-cycle-detecting-generational-garbage-collector}{%
\subsection{3.4 Python 2.0 Cycle-Detecting Generational Garbage
Collector}\label{python-2.0-cycle-detecting-generational-garbage-collector}}

Reference counting is deterministic and fast, but it cannot reclaim
self-referencing cycles:

\begin{lstlisting}[language=Python]
# Reference Cycle Example
a = []
b = []
a.append(b)
b.append(a)
del a, b  # reference counts of both objects remain 1, leaking memory!
\end{lstlisting}

To solve this, Python 2.0 introduced a generational cycle-detecting
Garbage Collector (GC) to periodically identify and sweep cycles.

\hypertarget{tracked-objects-and-gc-headers}{%
\subsubsection{1. Tracked Objects and GC
Headers}\label{tracked-objects-and-gc-headers}}

Only container objects (lists, dictionaries, tuples, custom classes) can
participate in reference cycles. Simple types like integers or strings
cannot. CPython prefixes all tracked container allocations with a
\passthrough{\lstinline!PyGC\_Head!} header, linking them into a
doubly-linked list.

\hypertarget{generational-strategy}{%
\subsubsection{2. Generational Strategy}\label{generational-strategy}}

The GC groups objects into three generations (Gen 0, Gen 1, Gen 2).
Newly allocated objects enter Gen 0. If they survive a garbage
collection pass, they are promoted to Gen 1, and eventually to Gen 2.
This strategy is based on the \textbf{weak generational hypothesis}:
most objects die young.

\hypertarget{cycle-isolation-algorithm}{%
\subsubsection{3. Cycle Isolation
Algorithm}\label{cycle-isolation-algorithm}}

\begin{enumerate}
\def\labelenumi{\arabic{enumi}.}
\tightlist
\item
  \textbf{Copy Reference Counts}: The GC copy-assigns the reference
  count of every object in the collection generation to its
  \passthrough{\lstinline!gc\_refs!} field in the
  \passthrough{\lstinline!PyGC\_Head!} header.
\item
  \textbf{Trace and Subtract}: The GC traverses the reference pointers
  of every object in the list. If a target object is also in the list,
  the GC decrements the target's \passthrough{\lstinline!gc\_refs!}
  value.
\item
  \textbf{Analyze Remainders}:

  \begin{itemize}
  \tightlist
  \item
    If an object's \passthrough{\lstinline!gc\_refs!} drops to zero, it
    is unreachable from outside the generation list (candidate for
    garbage).
  \item
    If \passthrough{\lstinline!gc\_refs > 0!}, the object is reachable
    from outside. The object and all objects it references are marked as
    reachable and promoted.
  \end{itemize}
\item
  \textbf{Deallocate}: The GC breaks cycles for remaining unreachable
  objects and frees their memory.
\end{enumerate}

\begin{center}\rule{0.5\linewidth}{0.5pt}\end{center}
